% Template for PLoS
% Version 3.4 January 2017
%
% % % % % % % % % % % % % % % % % % % % % %
%
% -- IMPORTANT NOTE
%
% This template contains comments intended 
% to minimize problems and delays during our production 
% process. Please follow the template instructions
% whenever possible.
%
% % % % % % % % % % % % % % % % % % % % % % % 
%
% Once your paper is accepted for publication, 
% PLEASE REMOVE ALL TRACKED CHANGES in this file 
% and leave only the final text of your manuscript. 
% PLOS recommends the use of latexdiff to track changes during review, as this will help to maintain a clean tex file.
% Visit https://www.ctan.org/pkg/latexdiff?lang=en for info or contact us at latex@plos.org.
%
%
% There are no restrictions on package use within the LaTeX files except that 
% no packages listed in the template may be deleted.
%
% Please do not include colors or graphics in the text.
%
% The manuscript LaTeX source should be contained within a single file (do not use \input, \externaldocument, or similar commands).
%
% % % % % % % % % % % % % % % % % % % % % % %
%
% -- FIGURES AND TABLES
%
% Please include tables/figure captions directly after the paragraph where they are first cited in the text.
%
% DO NOT INCLUDE GRAPHICS IN YOUR MANUSCRIPT
% - Figures should be uploaded separately from your manuscript file. 
% - Figures generated using LaTeX should be extracted and removed from the PDF before submission. 
% - Figures containing multiple panels/subfigures must be combined into one image file before submission.
% For figure citations, please use "Fig" instead of "Figure".
% See http://journals.plos.org/plosone/s/figures for PLOS figure guidelines.
%
% Tables should be cell-based and may not contain:
% - spacing/line breaks within cells to alter layout or alignment
% - do not nest tabular environments (no tabular environments within tabular environments)
% - no graphics or colored text (cell background color/shading OK)
% See http://journals.plos.org/plosone/s/tables for table guidelines.
%
% For tables that exceed the width of the text column, use the adjustwidth environment as illustrated in the example table in text below.
%
% % % % % % % % % % % % % % % % % % % % % % % %
%
% -- EQUATIONS, MATH SYMBOLS, SUBSCRIPTS, AND SUPERSCRIPTS
%
% IMPORTANT
% Below are a few tips to help format your equations and other special characters according to our specifications. For more tips to help reduce the possibility of formatting errors during conversion, please see our LaTeX guidelines at http://journals.plos.org/plosone/s/latex
%
% For inline equations, please be sure to include all portions of an equation in the math environment.  For example, x$^2$ is incorrect; this should be formatted as $x^2$ (or $\mathrm{x}^2$ if the romanized font is desired).
%
% Do not include text that is not math in the math environment. For example, CO2 should be written as CO\textsubscript{2} instead of CO$_2$.
%
% Please add line breaks to long display equations when possible in order to fit size of the column. 
%
% For inline equations, please do not include punctuation (commas, etc) within the math environment unless this is part of the equation.
%
% When adding superscript or subscripts outside of brackets/braces, please group using {}.  For example, change "[U(D,E,\gamma)]^2" to "{[U(D,E,\gamma)]}^2". 
%
% Do not use \cal for caligraphic font.  Instead, use \mathcal{}
%
% % % % % % % % % % % % % % % % % % % % % % % % 
%
% Please contact latex@plos.org with any questions.
%
% % % % % % % % % % % % % % % % % % % % % % % %

\documentclass[10pt,letterpaper]{article}
\usepackage[top=0.85in,left=2.75in,footskip=0.75in]{geometry}

% amsmath and amssymb packages, useful for mathematical formulas and symbols
\usepackage{amsmath,amssymb}

% Use adjustwidth environment to exceed column width (see example table in text)
\usepackage{changepage}

% Use Unicode characters when possible
\usepackage[utf8x]{inputenc}

% textcomp package and marvosym package for additional characters
\usepackage{textcomp,marvosym}

% cite package, to clean up citations in the main text. Do not remove.
\usepackage{cite}

% Use nameref to cite supporting information files (see Supporting Information section for more info)
\usepackage{nameref,hyperref}

% line numbers
\usepackage[right]{lineno}

% ligatures disabled
\usepackage{microtype}
\DisableLigatures[f]{encoding = *, family = * }

% color can be used to apply background shading to table cells only
\usepackage[table]{xcolor}

% array package and thick rules for tables
\usepackage{array}

% create "+" rule type for thick vertical lines
\newcolumntype{+}{!{\vrule width 2pt}}

% create \thickcline for thick horizontal lines of variable length
\newlength\savedwidth
\newcommand\thickcline[1]{%
  \noalign{\global\savedwidth\arrayrulewidth\global\arrayrulewidth 2pt}%
  \cline{#1}%
  \noalign{\vskip\arrayrulewidth}%
  \noalign{\global\arrayrulewidth\savedwidth}%
}

% \thickhline command for thick horizontal lines that span the table
\newcommand\thickhline{\noalign{\global\savedwidth\arrayrulewidth\global\arrayrulewidth 2pt}%
\hline
\noalign{\global\arrayrulewidth\savedwidth}}


% Remove comment for double spacing
%\usepackage{setspace} 
%\doublespacing

% Text layout
\raggedright
\setlength{\parindent}{0.5cm}
\textwidth 5.25in 
\textheight 8.75in

% Bold the 'Figure #' in the caption and separate it from the title/caption with a period
% Captions will be left justified
\usepackage[aboveskip=1pt,labelfont=bf,labelsep=period,justification=raggedright,singlelinecheck=off]{caption}
\renewcommand{\figurename}{Fig}

% Use the PLoS provided BiBTeX style
\bibliographystyle{plos2015}

% Remove brackets from numbering in List of References
\makeatletter
\renewcommand{\@biblabel}[1]{\quad#1.}
\makeatother

% Leave date blank
\date{}

% Header and Footer with logo
\usepackage{lastpage,fancyhdr,graphicx}
\usepackage{epstopdf}
\pagestyle{myheadings}
\pagestyle{fancy}
\fancyhf{}
\setlength{\headheight}{27.023pt}
\lhead{\includegraphics[width=2.0in]{PLOS-submission.eps}}
\rfoot{\thepage/\pageref{LastPage}}
\renewcommand{\footrule}{\hrule height 2pt \vspace{2mm}}
\fancyheadoffset[L]{2.25in}
\fancyfootoffset[L]{2.25in}
\lfoot{\sf PLOS}

%% Include all macros below

\newcommand{\lorem}{{\bf LOREM}}
\newcommand{\ipsum}{{\bf IPSUM}}

%% END MACROS SECTION


\begin{document}
\vspace*{0.2in}

% Title must be 250 characters or less.
\begin{flushleft}
{\Large
\textbf\newline{Amquery: a unified searchable database of 16S rRNA amplicon libraries} % Please use "sentence case" for title and headings (capitalize only the first word in a title (or heading), the first word in a subtitle (or subheading), and any proper nouns).
}
\newline
% Insert author names, affiliations and corresponding author email (do not include titles, positions, or degrees).
\\


Nikolay Romashchenko \textsuperscript{1 $\ast$},
Ilia Korvigo \textsuperscript{1, 2},
Evgeny Andronov \textsuperscript{1},

\bigskip
\textbf{1} Laboratory of Microbiological Monitoring and Bioremediation of Soils, All-Russia Research Institute for Agricultural Microbiology, St. Petersburg, Russia
\\
\textbf{2} Laboratory of Functional Genomics, Moscow Institute of Physics and Technology, Moscow, Russia
\\
\bigskip

% Insert additional author notes using the symbols described below. Insert symbol callouts after author names as necessary.
% 
% Remove or comment out the author notes below if they aren't used.
%
% Primary Equal Contribution Note
%\Yinyang These authors contributed equally to this work.

% Additional Equal Contribution Note
% Also use this double-dagger symbol for special authorship notes, such as senior authorship.
%\ddag These authors also contributed equally to this work.

% Current address notes
%\textcurrency Current Address: Dept/Program/Center, Institution Name, City, State, Country % change symbol to "\textcurrency a" if more than one current address note
% \textcurrency b Insert second current address 
% \textcurrency c Insert third current address

% Deceased author note
%\dag Deceased

% Group/Consortium Author Note
%\textpilcrow Membership list can be found in the Acknowledgments section.

% Use the asterisk to denote corresponding authorship and provide email address in note below.
* nikolay.romashchenko@gmail.com

\end{flushleft}
% Please keep the abstract below 300 words
\section*{Abstract}
Rapidly increasing amounts of 16S amplicon libraries require the development of efficient methods for similarity search. 
Here we present Amquery, a new tool for fast indexing and similarity search of 16S rRNA amplicon libraries, based on a technique of k-mer abundance comparison.
Experimental results on huge 16S dataset show that Amquery can significantly outperform clustering-based pipelines, implemented in tools such as QIIME, in the index construction and search tasks. 
Amquery is freely available at https://github.com/nromashchenko/amquery.

\linenumbers

% Use "Eq" instead of "Equation" for equation citations.
\section*{Introduction}
High-throughput sequencing of 16S amplicon libraries has become a very popular and powerful way of assessing the structure of various microbial communities \cite{qin2010human, seedorf2014bacteria, ligi2014characterization}.
As the number of published 16S amplicon libraries keeps growing so does the demand for a way to quickly process samples to find related communities similarly to the way BLAST for sequences. Nowadays a single amplicon library can easily contain many thousands of reads, making clustering-based pipelines, implemented in tools such as QIIME \cite{caporaso2010qiime}, inapplicable to large databases. To address these problems, we developed Amquery.

\pagebreak

\section*{Methods}
Amquery implements two major components: 1. a $k$-mer counting-based sample preprocessing routine and 2. a metric community dissimilarity measure.

\subsection*{k-mer counting and sample representation}
Given a word-length $k$, Amquery represents an arbitrary amplicon library as a discrete probability distribution over the corresponding $k$-mer space. To reduce memory consumption, the $k$-mers are encoded as quaternary values. More formally, each $k$-mer in a sequence $s = s_0,\dots,s_{n-1}$ is represented as its \textit{lexicographic rank} $\mathrm{lxr}_k(s)$, corresponding to the $k$-mer's position in the lexicographically sorted list of all possible $k$-mers.

\begin{eqnarray}
\label{eq:schemeP}
    \mathrm{lxr}_k(s) = \sum_{i=0}^{k−1} \mathrm{ord}(s_i) \cdot |{\mathcal{L}}|^{k−i−1}
\end{eqnarray}
Here $\mathrm{ord}(s_i)$ is a value of ordering function at $s_i$ over the alphabet ${\mathcal{L}} = \{A, C, G, T \}$. The most important advantage of this representation is that $\mathrm{lxr}_k$ is a special case of the Karp-Rabin hash fingerprint \cite{karp1987efficient}, making it possible to map $k$-mer $s_{i:i+k-1}$ to quaternary value $\mathrm{lxr}_k(s_{i:i+k-1})$, given the previous $k$-mer $s_{i−1:i+k-2}$ and its quaternary $\mathrm{lxr}_k(s_{i−1:i+k-2})$, in constant time:

\begin{eqnarray}
\label{eq:schemeP}
    \mathrm{lxr}_k(s_{i:i+k-1}) = |{\mathcal{L}}| \cdot (\mathrm{lxr}_k(s_{i−1:i+k-2}) − ord(s_{i−1}) \cdot 
									 |{\mathcal{L}}|^{k−1}) + ord(s_{i+k−1})
\end{eqnarray}
This advantage, however, imposes an upper-bound on the word-length $k$, since the number of possible $k$-mers grows exponentially with increasing word-length. That is, for a $k$-mer value to fit in a single machine word, we assume $k \leq 32$ for x64 architectures.
Since realistic sample ${k}$-mer distribution vectors are extremely high-dimensional and sparse, Amquery represents them as rank-sorted sparse arrays. In our implementation, a sparse array is a key-value storage in which the keys are the lexicographic ranks and the values are the empirical probabilities (i.e. relative abundances of the corresponding k-mers).

\subsection*{Distance metric}

After Amquery transforms a set of samples into discrete probability distributions, it is quite natural to use some metric on distributions to compare samples. Amquery utilises the square root of \textit{Jensen-Shannon divergence} (RJSD) \cite{lin1991divergence} that satisfies many statistical properties and fits many cases of high-dimensional spaces in $\mathbb{R^\star}$ as a distance metric \cite{fuglede2004jensen, endres2003new}. We make use of the specific form of RJSD, introduced by Connor et al. \cite{connor2013evaluation}, to speed up distance evaluation over sparse $k$-mer distributions.

\subsection*{Metric indexing}
Amquery utilises vantage-point tree (vp-tree) \cite{yianilos1993data, chavez2001searching} as an efficient data structure, providing fast metric-based indexing and search. Since vantage-point tree is designed to be independent of the space dimensionality \cite{yianilos1993data}, this allows us to store high-dimensional sparse vectors with virtually no loss of indexing and search performance: it takes $\mathcal{O}(n)$ space and $\mathcal{O}(n \log n)$ time in the worst case to be built (for an input set of $n$ samples), and a search query is argued to require $\mathcal{O}(\log n)$ of time, with some reservations. For more detailed analysis of vp-tree performance, see \cite{yianilos1993data}.

\section*{Results}
The implementation of Amquery was compared against our own custom solution of similarity search, based on the classic microbiome analysis QIIME pipeline, in the following fashion.
We computed pairwise weighted Unifrac distances \cite{lozupone2011unifrac} and Bray-Curtis similarities for huge 16S metagenomic dataset by performing both a reference-based and \textit{de novo} OTU picking.
Both of resulting pairwise distance matrices were used for naive top-k nearest neighbor selection as a reference to estimate search relevance of Amquery. 
We also compare speed and memory efficiency of index construction by Amquery against these QIIME based beta-diversity pipelines.
All of the experiments were run on a machine equipped with an Intel Xeon CPU E5-2650 v3 (2.30 GHz) and a Seagate Constellation ES.3 ST3000NM0033 hard drive.

\subsection*{Data}
We analyzed a human gut 16S metagenomic dataset, which consists of 1978 samples, obtained from the Sequence Read Archive \cite{leinonen2010sequence}. 
For each sample, the reads were filtered by length with 250 base pairs threshold and rarefied by choosing 1000 sequences per sample randomly.
Resulting samples are split into two parts: referred below to as \texttt{main}, the size of which varies in range (100, 200, $\dots$, 1000), and referred below to as \texttt{additional}, 
consist of all the rest samples, accordingly. Resulting pairs of (\texttt{main}, \texttt{additional}) sample subsets were used in further processing.

\subsection*{Data processing}
\subsubsection*{By Amquery}
For every split of the dataset, Amquery index was built for the \texttt{main} part, using $k=15$ as k-mer size. 100 random samples from the \texttt{main} part were searched in resulting index,
and a random subset of the \texttt{additional} sample set of size 100, 200, $\dots$, 1000 was added to the index separately. At every stage of processing, peak RAM usage and computational time were measured.

\subsubsection*{By QIIME}
We used QIIME versions 1.9.1 to perform OTU picking and beta-diversity analyses for the \texttt{main} part of every dataset split.
Reference-based and \textit{de novo} OTU picking were done using the default options as implemented in QIIME at the 95\% sequence similarity threshold, 
except for \texttt{uclust\_ref} OTU picking method for the reference-based case.

\subsection*{Speed and memory efficiency}
Elapsed time measurements at each stage of processing both for reference-based and \textit{de novo} QIIME pipeline against Amquery are presented in Fig~\ref{fig1} and Fig~\ref{fig2} respectively.
More detailed view of these measurements is presented Table~\ref{table1} and Table~\ref{table2},
where rows present measurements, made for each of the dataset splits, and the columns present elapsed time at each stage of pipelines. 
Similarly, peak RAM usage presented in Fig~\ref{fig3} and Fig~\ref{fig4} graphically and in Table~\ref{table3} and Table~\ref{table4} in numbers.


\begin{figure}[!h]
\caption{{\bf Elapsed time for QIIME (reference-based OTU picking) and Amquery pipelines for sets of samples of different size.}
A, B: Measurements of elapsed time, spended for reference-based OTU picking and beta-diversity analysis by QIIME using Bray-Curtis similarity (A) and weighted Unifrac (B). 
C: Measurements of elapsed time, spended for RJSD metric indexing by Amquery.}
\label{fig1}
\end{figure}

\begin{figure}[!h]
\caption{{\bf Elapsed time for QIIME (\textit{de novo} OTU picking) and Amquery pipelines for sets of samples of different size.}
A, B: Measurements of elapsed time, spended for \textit{de novo} OTU picking and beta-diversity analysis by QIIME using Bray-Curtis similarity (A) and weighted Unifrac (B). 
C: Measurements of elapsed time, spended for RJSD metric indexing by Amquery. Note that the C subfigure is the Fig~\ref{fig1}.C subfigure in the scale of A and B measurements. }
\label{fig2}
\end{figure}

\begin{table}[!ht]
\begin{adjustwidth}{-2.25in}{0in} % Comment out/remove adjustwidth environment if table fits in text column.
\centering
\caption{{\bf Measurements of processing time (mm:ss.ms) of Amquery and reference-based QIIME pipelines }}
\begin{tabular}{|l+l|l|l|l|l|l|l|l|l|}
\hline
\textit{size} & \multicolumn{3}{|l|}{\bf Amquery} & \multicolumn{6}{|l|}{\bf Reference-based QIIME}\\ \hline
& k-mers & indexing & {\bf Total} & pick\_otus & make\_otu & bdiv\_wu & bdiv\_bc & {\bf Total (wu)} & {\bf Total (bc)}  \\ \thickhline
100 & 0:43.08 & 0:00.74 & {\bf 0:43.82 } & 0:56.55 & 0:01.33 & 0:01.46 & 0:01.52 & {\bf 0:59.34} & {\bf 0:59.40} \\ \hline
200 & 1:25.16 & 0:01.51 & {\bf 1:26.67 } & 1:59.69 & 0:01.41 & 0:02.79 & 0:03.31 & {\bf 2:03.89} & {\bf 2:04.41} \\ \hline
300 & 2:10.24 & 0:02.40 & {\bf 2:12.64 } & 2:53.94 & 0:01.59 & 0:05.09 & 0:07.61 & {\bf 3:00.62} & {\bf 3:03.14} \\ \hline
400 & 2:44.32 & 0:03.13 & {\bf 2:47.45 } & 3:41.92 & 0:01.40 & 0:08.56 & 0:14.67 & {\bf 3:51.88} & {\bf 3:57.99} \\ \hline
500 & 3:37.39 & 0:04.31 & {\bf 3:41.70 } & 5:33.65 & 0:01.77 & 0:14.82 & 0:28.01 & {\bf 5:50.24} & {\bf 6:03.43} \\ \hline
600 & 4:17.47 & 0:05.45 & {\bf 4:22.92 } & 6:42.26 & 0:01.98 & 0:21.03 & 0:46.02 & {\bf 7:05.27} & {\bf 7:30.26} \\ \hline
700 & 5:08.55 & 0:06.00 & {\bf 5:14.55 } & 7:52.58 & 0:02.15 & 0:28.10 & 1:08.47 & {\bf 8:22.83} & {\bf 9:03.20} \\ \hline
800 & 5:38.59 & 0:07.07 & {\bf 5:45.66 } & 9:30.34 & 0:02.25 & 0:39.75 & 1:40.54 & {\bf 10:12.34} & {\bf 11:13.13} \\ \hline
900 & 6:33.67 & 0:07.83 & {\bf 6:41.50 } & 10:55.35 & 0:02.47 & 0:50.01 & 2:22.95 & {\bf 11:47.83} & {\bf 13:20.77} \\ \hline
1000 & 7:15.78 & 0:09.25 & {\bf 7:25.03 } & 11:54.37 & 0:02.54 & 1:04.56 & 2:58.86 & {\bf 13:01.47} & {\bf 14:55.77} \\ \hline
\end{tabular}
\begin{flushleft}
Interpretation of headings: \textit{size} is the number of samples, used for index construction. 
\textit{k-mers} is the time spent on k-mer counting.
\textit{indexing} is the time spent on vp-tree construction, including calculation of necessary RJSD distance values.
\textit{Total} is the total time of processing by Amquery.
\textit{pick\_otus} and \textit{make\_otu} are the total time spent by the \texttt{pick\_otus.py} and the \texttt{make\_otu\_table.py} scripts accordingly.
\textit{bdiv\_wu} and \textit{bdiv\_bc} are the total time spent by the \texttt{beta\_diversity.py} script for weighted Unifrac and Bray-Curtis similarities accordingly.
\textit{Total (wu)} and \textit{Total (bc)} are the total time spent by the QIIME pipeline for weighted Unifrac and Bray-Curtis similarities accordingly.
\end{flushleft}
\label{table1}
\end{adjustwidth}
\end{table}



\begin{table}[!ht]
\begin{adjustwidth}{-2.25in}{0in} % Comment out/remove adjustwidth environment if table fits in text column.
\centering
\caption{{\bf Measurements of processing time (mm:ss.ms or hh:mm:ss.ms) of \textit{de novo} based QIIME pipeline }}
\begin{tabular}{|l+l|l|l|l|l|l|l|l|l|l|}
\hline
\textit{size} & \multicolumn{10}{|l|}{\bf \textit{De novo} based QIIME}\\ \hline
& pick\_otus & rep\_set & align & filt\_aln & make\_phyl & make\_otu & bdiv\_wu & bdiv\_bc & {\bf Total (wu)} & {\bf Total (bc)}  \\ \thickhline
100 & 0:30.08 & 0:02.73 & 8:16.56 & 0:26.63 & 2:48.41 & 0:01.29 & 0:02.48 & 0:01.28 & {\bf 12:08.18} & {\bf 12:06.98} \\ \hline
200 & 1:03.83 & 0:04.11 & 15:33.94 & 0:46.64 & 5:10.33 & 0:01.46 & 0:07.04 & 0:03.14 & {\bf 22:47.35} & {\bf 22:43.45} \\ \hline
300 & 1:48.02 & 0:05.69 & 21:56.02 & 1:15.08 & 8:42.55 & 0:01.38 & 0:21.08 & 0:06.89 & {\bf 34:09.82} & {\bf 33:55.63} \\ \hline
400 & 2:25.44 & 0:06.97 & 26:42.65 & 1:26.78 & 11:38.27 & 0:01.52 & 0:40.01 & 0:13.04 & {\bf 43:01.64} & {\bf 42:34.67} \\ \hline
500 & 3:33.26 & 0:08.66 & 36:45.55 & 2:01.01 & 17:12.79 & 0:01.82 & 1:22.73 & 0:25.17 & {\bf 01:01:05.82} & {\bf 01:00:08.26} \\ \hline
600 & 4:38.02 & 0:10.59 & 42:20.45 & 2:26.52 & 22:15.39 & 0:01.75 & 2:22.34 & 0:42.02 & {\bf 01:14:15.06} & {\bf 01:12:34.74} \\ \hline
700 & 5:35.56 & 0:11.72 & 48:38.01 & 2:37.39 & 24:16.03 & 0:01.90 & 3:28.34 & 1:04.51 & {\bf 01:24:48.95} & {\bf 01:22:25.12} \\ \hline
800 & 7:03.39 & 0:13.10 & 55:14.31 & 2:55.26 & 27:43.87 & 0:02.19 & 7:49.36 & 1:32.94 & {\bf 01:41:01.48} & {\bf 01:34:45.06} \\ \hline
900 & 8:16.08 & 0:14.90 & 1:00:22 & 3:32.20 & 32:22.70 & 0:02.25 & 7:31.52 & 2:09.78 & {\bf 01:52:21.65} & {\bf 01:46:59.91} \\ \hline
1000 & 9:08.44 & 0:16.10 & 1:01:21 & 3:42.98 & 36:39.14 & 0:02.41 & 9:46.51 & 2:52.90 & {\bf 02:00:56.58} & {\bf 01:54:02.97} \\ \hline
\end{tabular}
\begin{flushleft}

\textit{pick\_otus}, \textit{make\_otu}, \textit{bdiv\_wu}, \textit{bdiv\_bc}, \textit{Total (wu)} and \textit{Total (bc)} columns have the same meanings as in the Table~\ref{table1}.
\textit{rep\_set}, \textit{align}, \textit{filt\_aln} and \textit{make\_phyl} contain the total time spent by the
\texttt{pick\_rep\_set.py}, \texttt{align\_seqs.py}, \texttt{filter\_alignment.py} and \texttt{make\_phylogeny.py} scripts accordingly.

\end{flushleft}
\label{table2}
\end{adjustwidth}
\end{table}



\begin{figure}[!h]
\caption{{\bf Peak RAM Usage for QIIME (reference-based OTU picking) and Amquery pipelines for sets of samples of different size.}
A, B: Peak RAM usage during reference-based OTU picking and beta-diversity analysis by QIIME using Bray-Curtis similarity (A) and weighted Unifrac (B). 
C: Peak RAM usage during sample indexing by Amquery.}
\label{fig3}
\end{figure}

\begin{figure}[!h]
\caption{{\bf Peak RAM Usage for QIIME (\textit{de novo} OTU picking) and Amquery pipelines for sets of samples of different size.}
A, B: Peak RAM usage during \textit{de novo} OTU picking and beta-diversity analysis by QIIME using Bray-Curtis similarity (A) and weighted Unifrac (B). 
C: Peak RAM usage during sample indexing by Amquery, the same as Fig~\ref{fig3}.C subfigure in the scale of A and B measurements.}
\label{fig4}
\end{figure}



\begin{table}[!ht]
\begin{adjustwidth}{-2.25in}{0in} % Comment out/remove adjustwidth environment if table fits in text column.
\centering
\caption{{\bf Peak RAM usage of Amquery and reference-based QIIME pipelines }}
\begin{tabular}{|l+l|l|l|l|l|l|l|}
\hline
\textit{size} & {\bf Amquery} & \multicolumn{6}{|l|}{\bf Reference-based QIIME}\\ \hline
& {\bf Peak} & pick\_otus & make\_otu & bdiv\_wu & bdiv\_bc & {\bf Peak (wu)} & {\bf Peak (bc)}  \\ \thickhline
100 & 100Mb & 242Mb & 120Mb & 145Mb & 135Mb & 242Mb & 242Mb \\ \hline
200 & 149Mb & 344Mb & 130Mb & 209Mb & 197Mb & 344Mb & 344Mb \\ \hline
300 & 222Mb & 437Mb & 135Mb & 295Mb & 275Mb & 437Mb & 437Mb \\ \hline
400 & 229Mb & 499Mb & 139Mb & 387Mb & 366Mb & 499Mb & 499Mb \\ \hline
500 & 290Mb & 649Mb & 146Mb & 541Mb & 516Mb & 649Mb & 649Mb \\ \hline
600 & 334Mb & 759Mb & 153Mb & 716Mb & 686Mb & 759Mb & 759Mb \\ \hline
700 & 379Mb & 833Mb & 163Mb & 881Mb & 850Mb & 881Mb & 850Mb \\ \hline
800 & 428Mb & 911Mb & 168Mb & 1.0Gb & 1.0Gb & 1.0Gb & 1.0Gb \\ \hline
900 & 474Mb & 995Mb & 175Mb & 1.3Gb & 1.2Gb & 1.3Gb & 1.2Gb \\ \hline
1000 & 515Mb & 1.0Gb & 183Mb & 1.4Gb & 1.4Gb & 1.4Gb & 1.4Gb \\ \hline
\end{tabular}
\begin{flushleft}
Column naming notation corresponds to the Table~\ref{table1}. \textit{Peak} columns present the peak RAM usage of corresponding pipelines.   
\end{flushleft}
\label{table3}
\end{adjustwidth}
\end{table}



\begin{table}[!ht]
\begin{adjustwidth}{-2.25in}{0in} % Comment out/remove adjustwidth environment if table fits in text column.
\centering
\caption{{\bf Peak RAM usage of \textit{de novo} based QIIME pipeline }}
\begin{tabular}{|l+l|l|l|l|l|l|l|l|l|l|}
\hline
\textit{size} & \multicolumn{10}{|l|}{\bf \textit{De novo} based QIIME}\\ \hline
& pick\_otus & rep\_set & align & filt\_aln & make\_phyl & make\_otu & bdiv\_wu & bdiv\_bc & {\bf Peak (wu)} & {\bf Peak (bc)}  \\ \thickhline

100 & 174Mb & 174Mb & 298Mb & 166Mb & 149Mb & 118Mb & 169Mb & 132Mb & 298Mb & 298Mb \\ \hline
200 & 250Mb & 234Mb & 424Mb & 211Mb & 173Mb & 123Mb & 259Mb & 184Mb & 424Mb & 424Mb \\ \hline
300 & 345Mb & 295Mb & 578Mb & 269Mb & 205Mb & 131Mb & 407Mb & 264Mb & 578Mb & 578Mb \\ \hline
400 & 404Mb & 356Mb & 638Mb & 291Mb & 281Mb & 137Mb & 541Mb & 353Mb & 638Mb & 638Mb \\ \hline
500 & 545Mb & 420Mb & 822Mb & 363Mb & 377Mb & 142Mb & 793Mb & 492Mb & 822Mb & 822Mb \\ \hline
600 & 645Mb & 476Mb & 976Mb & 422Mb & 462Mb & 151Mb & 1Gb & 650Mb & 1Gb & 976Mb\\ \hline
700 & 716Mb & 553Mb & 1.0Gb & 445Mb & 500Mb & 161Mb & 1.3Gb & 803Mb & 1.3Gb & 1.0Gb\\ \hline
800 & 791Mb & 607Mb & 1.1Gb & 478Mb & 548Mb & 167Mb & 1.6Gb & 981Mb & 1.6Gb & 1.1Gb\\ \hline
900 & 871Mb & 666Mb & 1.3Gb & 548Mb & 570Mb & 174Mb & 2.Gb & 1.1Gb & 2.Gb & 1.3Gb\\ \hline
1000 & 903Mb & 726Mb & 1.4Gb & 568Mb & 676Mb & 179Mb & 2.2Gb & 1.3Gb & 2.2Gb & 1.4Gb \\ \hline

\end{tabular}
\begin{flushleft}
Column namings correspond to the tables above.   
\end{flushleft}
\label{table4}
\end{adjustwidth}
\end{table}




\subsection*{Similarity search effectiveness}
To estimate relevance of similarity search, provided by Amquery, mean precision at k (MP@K), mean average precision at k (MAP@K), and normalized cumulative discounted gain at k (NCDG@k) \cite{book, jarvelin2000ir, jarvelin2002cumulated}
are evaluated, with the relevance function, declared as follows. For every search query $q$, all the pairwise distances $D(q, x)$ provided by QIIME pipeline from $q$ to every sample $x$ in the index, are normalized. 
Thus, resulting relevance value $R(q, r)$ for search output result $r$ is the value $1 - \hat{D}(q, r)$, whree $\hat{D}$ is the normalized value of $D(q, x)$ over all the $x$ in the index. Here we use the notation of $D(q, r)$ for weighted Unifrac and Bray-Curtis similarity distances.

All the metrics are evaluated among all the dataset splits, mentioned above, for search query size $k = {1, 3, 5, 7, 10, 15, 20}$ and 100 random samples from the index. Since precision metrics as MP@k and MAP@k require binary relevance function,
$R(q, r)$ values for these metrics are binarized as follows:

\[ R_b(q, r)  =
  \begin{cases}
    1  & r \text{ in top-} k \text{ nearest samples of } q, \text{ according to } D(q, r)\\
    0  &  otherwise \\
  \end{cases}
\] 

MP@k, MAP@k, NDCG@k evaluation results presented in the Fig~\ref{fig5}.


\begin{figure}[!h]
\caption{{\bf MP@k (A, D), MAP@k (B, E) and NDCG@k (C, F) for Amquery search queries and random sampling from the index as a baseline, according to relevance functions, based on the reference-based QIIME pipeline result.}
A, B, C are calculated according to weighted Unifrac, while D, E, F are calculated according to Bray-Curtis similarity as a relevance function.}
\label{fig5}
\end{figure}


\subsection*{Conclusions}
We presented a k-mer based approach for fast indexing and searching of 16S amplicon libraries, which designed to outperform clustering-based solutions for similarity search, in terms of speed and memory consumption.
This approach allows for indexing thousands of samples in minutes, and lightning fast similarity search. We demonstrated that similarity search results, based on RJSD distance over estimated k-mer distributions, correspond to a certain extent to the weighted Unifrac based and Bray-Curtis similarity based search results.


\begin{thebibliography}{10}

\bibitem{caporaso2010qiime}
J~Gregory Caporaso, Justin Kuczynski, Jesse Stombaugh, Kyle Bittinger,
  Frederic~D Bushman, Elizabeth~K Costello, Noah Fierer, Antonio~Gonzalez
  Pe{\~n}a, Julia~K Goodrich, Jeffrey~I Gordon, et~al.
\newblock Qiime allows analysis of high-throughput community sequencing data.
\newblock {\em Nature methods}, 7(5):335--336, 2010.

\bibitem{chavez2001searching}
Edgar Ch{\'a}vez, Gonzalo Navarro, Ricardo Baeza-Yates, and Jos{\'e}~Luis
  Marroqu{\'\i}n.
\newblock Searching in metric spaces.
\newblock {\em ACM computing surveys (CSUR)}, 33(3):273--321, 2001.

\bibitem{book}
Prabhakar~Raghavan Christopher D~Manning and Hinrich Schütze.
\newblock {\em Introduction to Information Retrieval}.
\newblock Cambridge University Press, 2008.

\bibitem{connor2013evaluation}
Richard Connor, Franco~Alberto Cardillo, Robert Moss, and Fausto Rabitti.
\newblock Evaluation of jensen-shannon distance over sparse data.
\newblock In {\em International Conference on Similarity Search and
  Applications}, pages 163--168. Springer, 2013.

\bibitem{endres2003new}
Dominik~Maria Endres and Johannes~E Schindelin.
\newblock A new metric for probability distributions.
\newblock {\em IEEE Transactions on Information theory}, 49(7):1858--1860,
  2003.

\bibitem{fuglede2004jensen}
Bent Fuglede and Flemming Topsoe.
\newblock Jensen-shannon divergence and hilbert space embedding.
\newblock In {\em Information Theory, 2004. ISIT 2004. Proceedings.
  International Symposium on}, page~31. IEEE, 2004.

\bibitem{jarvelin2000ir}
Kalervo J{\"a}rvelin and Jaana Kek{\"a}l{\"a}inen.
\newblock Ir evaluation methods for retrieving highly relevant documents.
\newblock In {\em Proceedings of the 23rd annual international ACM SIGIR
  conference on Research and development in information retrieval}, pages
  41--48. ACM, 2000.

\bibitem{jarvelin2002cumulated}
Kalervo J{\"a}rvelin and Jaana Kek{\"a}l{\"a}inen.
\newblock Cumulated gain-based evaluation of ir techniques.
\newblock {\em ACM Transactions on Information Systems (TOIS)}, 20(4):422--446,
  2002.

\bibitem{karp1987efficient}
Richard~M Karp and Michael~O Rabin.
\newblock Efficient randomized pattern-matching algorithms.
\newblock {\em IBM Journal of Research and Development}, 31(2):249--260, 1987.

\bibitem{leinonen2010sequence}
Rasko Leinonen, Hideaki Sugawara, and Martin Shumway.
\newblock The sequence read archive.
\newblock {\em Nucleic acids research}, page gkq1019, 2010.

\bibitem{ligi2014characterization}
Teele Ligi, Kristjan Oopkaup, Marika Truu, Jens-Konrad Preem, Hiie N{\~o}lvak,
  William~J Mitsch, {\"U}lo Mander, and Jaak Truu.
\newblock Characterization of bacterial communities in soil and sediment of a
  created riverine wetland complex using high-throughput 16s rrna amplicon
  sequencing.
\newblock {\em Ecological Engineering}, 72:56--66, 2014.

\bibitem{lin1991divergence}
Jianhua Lin.
\newblock Divergence measures based on the shannon entropy.
\newblock {\em IEEE Transactions on Information theory}, 37(1):145--151, 1991.

\bibitem{lozupone2011unifrac}
Catherine Lozupone, Manuel~E Lladser, Dan Knights, Jesse Stombaugh, and Rob
  Knight.
\newblock Unifrac: an effective distance metric for microbial community
  comparison.
\newblock {\em The ISME journal}, 5(2):169, 2011.

\bibitem{qin2010human}
Junjie Qin, Ruiqiang Li, Jeroen Raes, Manimozhiyan Arumugam,
  Kristoffer~Solvsten Burgdorf, Chaysavanh Manichanh, Trine Nielsen, Nicolas
  Pons, Florence Levenez, Takuji Yamada, et~al.
\newblock A human gut microbial gene catalogue established by metagenomic
  sequencing.
\newblock {\em nature}, 464(7285):59--65, 2010.

\bibitem{seedorf2014bacteria}
Henning Seedorf, Nicholas~W Griffin, Vanessa~K Ridaura, Alejandro Reyes, Jiye
  Cheng, Federico~E Rey, Michelle~I Smith, Gabriel~M Simon, Rudolf~H
  Scheffrahn, Dagmar Woebken, et~al.
\newblock Bacteria from diverse habitats colonize and compete in the mouse gut.
\newblock {\em Cell}, 159(2):253--266, 2014.

\bibitem{yianilos1993data}
Peter~N Yianilos.
\newblock Data structures and algorithms for nearest neighbor search in general
  metric spaces.
\newblock In {\em SODA}, volume~93, pages 311--21, 1993.

\end{thebibliography}



\end{document}

